\section{结论}
\label{sec:conclusion}

我们将图检查与利用源文本和字段约束的修复结合。评估概况记录局部、图和来源错误。主路径先清理输入，再请求模型生成完整图，最后过滤三元组。尺度感知选择器还在图副本上比较编辑，考虑质量增益、成本和图边界。共享提示以及过滤前后的评分展示各步骤带来的变化。

诊断＋门控 修复了 98.00\% 的注入缺陷，并将模型构建图的 F1 从 87.99\% 提高到 92.48\%。直接字段复制恢复了中文记录中的全部参考。来源行提示使收据 F1 比 简单流程 高 4.76 个百分点。后续测试中，完整索引与仅锚点的平均 F1 相同。合同索引修复的 F1 为 44.84\%，输入为 46.76\%；与重新抽取相差 0.72 个百分点（95\% 置信区间：$[-4.80,6.48]$）。来源和文档编号检查去掉了不需要的三元组。匹配测试中，错误报告降低了平均 F1；模型构建输入上，默认学习选择器的分数低于等权重设置。

人工评价在 95/200 项参考差异中发现输入错误。评审接受五种设置中的 134/199 项编辑，另有一项不确定。这些结果给出清楚的改进目标：为每个字段选择正确取值，保留正确的输入事实，并检查最终图。